\chapter*{De ce}
\addcontentsline{toc}{chapter}{De ce}

\begin{quote}
Toate modelele sunt gresite, dar unele sunt utile.

\begin{flushright}
-- George Box
\end{flushright}
\end{quote}

Inca ezit putin sa ma apuc de predat \textit{AI} pentru olimpiada:

Nu sunt foarte convins ca numele \textit{AI} descrie ceea ce facem prin
aceste metode. Mi se pare, de multe ori, o umbra palida a cognitiei: luam
lucruri extrem de complexe si le trecem printr-un model surogat si miop, optimizat
pentru un proxy care nu surprinde neaparat ceea ce ne intereseaza cu
adevarat. Iar termenul este folosit foarte des cu o conotatie mult mai
pretentioasa decat continutul real.

Mai important, cred ca o olimpiada de ML risca sa puna prea mult accent pe
\textbf{cum facem ceva} si prea putin pe \textbf{de ce facem acel lucru}.

La informatica, o metrica este de obicei destul de cinstita: o solutie
corecta este corecta. In ML, lucrurile sunt mai ample. Poti ajunge sa
optimizezi o metrica care este doar un proxy pentru un obiectiv mult mai
complex, iar raspunsul la o problema poate deveni pur si simplu ``facem o
retea mai mare''. In industrie se ajunge uneori sa construim sisteme care
optimizeaza foarte bine lucruri care poate nu ar fi trebuit optimizate in
primul rand.

De asta vreau ca acest curs sa fie mai mult decat o colectie de algoritmi
si trucuri pentru olimpiada.

\textbf{Magia apare atunci cand intelegi ce poti sa rezolvi, cum poti sa
rezolvi si, mai ales, de ce merita sa rezolvi problema respectiva.}

\section*{Creierul nu e facut pentru statistica}

Traim in cea mai sigura perioada din istoria omenirii, si totusi ne e
frica de aproape orice. Motivul nu e ca lumea reala e periculoasa, e ca
noi, ca specie, suntem prost echipati sa intelegem probabilitatile.
Creierul judeca riscul dupa cat de vie e o poveste, nu dupa cat de des se
intampla ceva de fapt. Ne e frica de avion, unde riscul de deces e undeva
la 1 la cateva milioane de zboruri, dar urcam linistiti la volan, unde
riscul e cu mult mai mare, doar ca nu are aceeasi incarcatura emotionala.

Acelasi mecanism il vezi la loterie: oamenii aleg numere ``norocoase'' sau
evita 13, desi fiecare combinatie are exact aceeasi sansa. Il vezi la
ruleta, cand dupa opt rosii la rand toata lumea e convinsa ca urmeaza
negru, desi bila nu tine minte nimic din rotirile anterioare (asta se
cheama \textit{falacia jucatorului}). Si il vezi la poker: un jucator bun
castiga pe termen lung nu pentru ca ghiceste cartile, ci pentru ca ia,
mana dupa mana, decizii corecte statistic, chiar daca pierde uneori o
mana anume. Diferenta dintre un jucator bun si unul prost nu e norocul de
moment, e intelegerea probabilitatilor pe termen lung, si o sa exersam
exact chestia asta la clasa, cu carti pe masa, nu doar cu formule pe
tabla.

De ce conteaza asta pentru un curs de ML? Pentru ca jobul vostru, ca
oameni care lucreaza cu date, e exact opusul acestui reflex: sa scoateti
decizia din mana intuitiei si sa o puneti pe o baza masurabila. Dar ca sa
faceti asta bine, trebuie sa fiti foarte atenti de unde vin datele alea.

\section*{Garbage in, garbage out}

Poate ati observat ca studiile de nutritie se contrazic constant: intr-un
an oule sunt periculoase, in urmatorul sunt super-aliment. Motivul, de
multe ori, nu e ca stiinta se schimba, e ca datele sunt proaste de la
inceput. Multe studii se bazeaza pe chestionare in care oamenii isi
amintesc ce au mancat saptamana trecuta, memorie nesigura, si raporteaza
ce cred ca ar trebui sa fi mancat, nu ce au mancat cu adevarat.

Nu conteaza cat de sofisticat e modelul vostru: o retea neuronala cu sute
de milioane de parametri, antrenata pe date subiective, tot zgomot invata,
nu realitate. Algoritmul e o oglinda, reflecta exact ce ii dai, buguri si
minciuni incluse. De-aia o sa insistam, in cursurile urmatoare, pe cum se
colecteaza si se curata datele, nu doar pe cum se antreneaza un model.

Un exemplu istoric bun pentru genul asta de eroare, dar dintr-un unghi
diferit, e \textit{Clever Hans}, un cal din Germania de la inceputul
secolului XX care parea sa stie sa faca operatii matematice, batand din
copita de cate ori era nevoie. Nu stia matematica, desigur: citea, fara sa
vrea nimeni asta, semnale corporale involuntare de la antrenorul lui, care
se relaxa vizibil cand calul ajungea la numarul corect. Asta se cheama,
azi, \textit{data leakage} sau \textit{shortcut learning}: modelul, sau calul, invata sa exploateze un
semnal care e corelat cu raspunsul corect, dar care n-are nicio legatura
reala cu problema, si care de multe ori nici n-o sa mai fie disponibil
cand rulezi modelul in productie. Ideea merita retinuta devreme: cand un
model merge suspect de bine, prima intrebare nu e ``ce genial e'', e ``ce
a gasit, de fapt, sa trisheze''.

\section*{Corelatie nu inseamna cauzalitate}

Japonezii mananca putina carne si traiesc foarte mult. Spaniolii mananca
de doua ori mai multa carne, de patru ori mai multa grasime animala, si
totusi sunt pe locul doi in lume la speranta de viata, la o distanta
infima de Japonia. Daca ai antrena un model simplu care leaga direct
consumul de carne de mortalitate, ai rata complet explicatia.

Ce lipseste din poza e tot ce nu ai masurat: sistemul de sanatate, cat
mers pe jos face fiecare zi, cat de relaxat e contextul social in care
mananci. Astea se numesc \textit{variabile de confuzie}, factori ascunsi
care influenteaza si cauza aparenta si efectul, dar pe care nu i-ai
inclus in date. Un model nu stie ce ai uitat sa masori, el vede doar ce
ii dai, si gaseste corelatii, nu explicatii. Partea grea a lucrului cu
date nu e sa rulezi un algoritm, e sa te gandesti ce lipseste din tabelul
tau inainte sa tragi o concluzie.

\section*{Metrica gresita, decizie gresita}

Un exemplu la scara mare: dupa accidentul de la Fukushima, unde numarul
de decese directe cauzate de radiatii a fost aproape zero, mai multe tari
au inchis centrale nucleare functionale, de frica unui accident extrem de
rar, si au trecut, partial, pe carbune. Poluarea din arderea carbunelui
insa e o cauza constanta si masurabila de mortalitate, an de an, la o
scara mult mai mare decat orice accident nuclear din istorie. Optimizand
impotriva riscului mai spectaculos, dar mai rar, ai putea agrava riscul
mai banal, dar mai frecvent si mai mare.

Asta e exact problema unei functii de cost (\textit{loss function}) prost
alese intr-un model de ML: daca optimizezi pentru metrica gresita, modelul
tau o sa ``rezolve'' cu succes o problema care nu era cea reala. O sa
vorbim mult despre metrici mai incolo in curs, accuracy, precision, recall
si restul, dar tineti minte de pe acum: alegerea metricii nu e un detaliu
tehnic, e decizia care conteaza cel mai mult.

Exista si o lege scurta care rezuma toata problema asta: cand o masura
devine tinta, ea inceteaza sa mai fie o masura buna, \textit{legea lui
Goodhart}, formulata ulterior si de \textit{Campbell}: cu cat un indicator
cantitativ social este folosit mai mult pentru luarea deciziilor, cu atat
va fi mai supus presiunilor de coruptie si va distorsiona procesul pe care
intentioneaza sa il monitorizeze. O vezi peste tot: note care devin scopul scolii, in loc de
invatare; KPI-uri de companie care se optimizeaza fara sa mai reflecte
performanta reala; retele sociale care optimizeaza timpul petrecut pe
aplicatie, nu binele utilizatorului. Fiecare model de ML pe care il
construiti o sa optimizeze exact ce ii cereti, nici un pic mai mult si
nici un pic mai intelept. Responsabilitatea sa cereti lucrul corect
ramane a voastra.

\section*{Cum o sa invatam}

O sa incepem cu ceva simplu si contraintuitiv, ca sa vedeti cat de prost
va poate insela intuitia: problema Monty Hall, cu cele trei usi, unde
raspunsul corect nu e cel pe care il da aproape toata lumea din prima. E
un exercitiu de cinci minute care merita mai mult decat o ora de teorie.

N-o sa fie doar teorie si cod pe ecran. O sa incerc sa va invat meseria
asta prin exercitii practice, si multe dintre ele o sa fie imprumutate
din jocuri: o sa jucam barbut, o sa jucam poker, si o sa calculam
impreuna sansele reale din spatele fiecarei decizii. Un joc de noroc e,
de fapt, un laborator perfect de probabilitati: regulile sunt simple,
miza e clara, si poti sa masori exact cat de bine gandesti fata de cat de
bine crezi ca gandesti.

O sa vorbim si despre lucruri din viata de zi cu zi, nu doar despre
exemple abstracte de manual. Vreau sa incepeti sa observati cat de des va
loviti de fapt de rationament probabilistic, chiar si atunci cand nu il
numiti asa: cand decideti daca luati umbrela, cand cantariti un risc, cand
cititi o stire cu un procent aruncat fara context.

O sa atingem si zone mai fine de optimizare matematica, dar mereu cu
intrebarea in minte: cat de des ne lovim de fapt de genul asta de
problema? Nu vreau sa invatati o formula doar ca sa o uitati dupa
concurs, vreau sa recunoasteti situatia in care chiar merita folosita.

Un obiectiv mai ambitios: la finalul cursului, sper sa puteti purta un
debate corect cu cineva, pe baza de date si probabilitati, nu pe baza de
cine vorbeste mai tare sau cine are povestea mai convingatoare. Asta e,
pana la urma, toata ideea: sa inlocuiti argumentul emotional cu unul care
rezista la verificare.

\section*{O nota de filozofie}

Undeva la baza tuturor acestora sta o intrebare mai veche decat
calculatoarele: cum ar trebui sa gandim in conditii de incertitudine?
Filozofii au numit-o \textit{problema inductiei} cu multe secole inainte
sa existe Machine Learning: din faptul ca soarele a rasarit in fiecare zi
din viata ta, nu poti fi sigur 100\% ca rasare si maine, doar foarte,
foarte increzator. Toata statistica moderna, si tot ce construim noi in
cursul asta, e un raspuns practic la intrebarea aia: cum cuantificam
increderea, in loc sa o lasam la nivel de senzatie.

Nu e o intrebare doar tehnica. E, in fond, o intrebare despre onestitate
intelectuala: esti dispus sa iti schimbi parerea cand apar date noi, sau
te agati de prima poveste care ti s-a parut convingatoare? Un model bun
de ML face asta automat, isi actualizeaza increderea pe masura ce vede
date noi. Un om rational ar trebui sa faca la fel.

Multe dintre dezbaterile de azi despre inteligenta artificiala si
transumanism nu sunt, de fapt, chiar noi. Intrebarea daca o masina poate
sa ``gandeasca'' cu adevarat, daca poate avea constiinta, daca deciziile
ei sunt cu adevarat ale ei sau doar rezultatul mecanic al unor cauze
anterioare, sunt variante reimpachetate ale unor dezbateri vechi de mii
de ani: ce e sufletul, exista liber arbitru sau totul e determinism, ce
ne face de fapt oameni. Filozofii dezbateau determinismul cu mult inainte
sa existe retele neuronale, teologii dezbateau natura sufletului cu mult
inainte sa existe calculatoare. Tehnologia noua ne da un vocabular nou
pentru intrebari vechi, nu neaparat raspunsuri noi la ele.

N-o sa rezolvam aici, in doua paragrafe, ce n-au rezolvat filozofii in
doua mii de ani. Dar data viitoare cand cineva va vinde o poveste
dramatica despre o inteligenta artificiala care ne va inlocui sau despre
masini constiente, intrebati-va cate din argumentele alea sunt de fapt
argumente teologice vechi, rescrise cu cuvinte noi.

Asta nu inseamna sa subestimati unealta: ML rezolva azi probleme pe care
acum douazeci de ani le credeam imposibile, si ne ajuta sa evoluam, nu
doar sa fim mai eficienti. Nu vreau sa iesiti din curs nici
naiv-entuziasmati, nici speriati de ziua de maine, ci suficient de
lamuriti incat sa folositi unealta cu incredere, nu cu superstitie.

\section*{A patra lovitura}

Freud numea asta cele trei rani ale orgoliului uman. Copernic ne-a scos
din centrul spatiului - Pamantul e o piatra oarecare, in jurul unei
stele oarecare. Darwin ne-a scos din centrul vietii - suntem o ramura
printre altele, nu o creatie separata. Iar psihanaliza, spunea chiar
Freud despre propria lui munca, ne-a scos din centrul mintii noastre: nu
esti stapan nici macar pe gandurile tale, pentru ca o mare parte din
ceea ce te face sa gandesti, sa simti si sa alegi se intampla fara tine.

Neurostiinta a dus rana mai departe. Robert Sapolsky merge pana la
capatul incomod al acestei idei: ceea ce numim o alegere este rezultatul
unei istorii care a inceput inainte ca noi sa putem alege - gene,
hormoni, creier, copilarie, mediu, cultura, tot ce ni s-a intamplat si
tot ce am mostenit. Nu trebuie sa acceptam concluzia lui despre liberul
arbitru ca fiind definitiva ca sa simtim forta intrebarii.
\textit{Cat din ceea ce numesc ``eu'' a fost, de fapt, vreodata al meu?}

Dar e o diferenta intre primele trei si ceea ce vine acum. Copernic,
Darwin si Freud au descoperit ceva ce era deja acolo. Lumea era asa
inainte sa o intelegem noi. Pamantul nu era in centru si inainte sa
aflam asta. Evolutia se intampla si inainte sa-i dam un nume. Iar faptul
ca mintea noastra nu este transparenta pentru noi nu a inceput odata cu
Freud. Noi doar am ajuns, treptat, sa vedem.

\textbf{AI e diferit. AI nu vine din afara povestii. Il construim noi.}

Nu descoperim ca inteligenta n-a fost niciodata exclusiv a noastra.
Facem, cu mana noastra, ceva care incepe sa ne atinga exact in locul in
care ne credeam de neinlocuit. De-asta doare altfel. Primele trei
lovituri ne-au facut mai putin centrali. AI ne poate face sa ne intrebam
daca mai suntem necesari.

Si poate nici macar nu trebuie sa ajungem la ideea unei civilizatii
artificiale ca sa simtim asta. Nu trebuie sa existe un oras populat de
fiinte digitale care se iubesc, se cearta si isi fac planuri. E
suficient un elev care, la un moment dat, se uita la o tema si se
gandeste:

\textit{``Ce sens are sa invat asta daca AI-ul stie deja?''}

Aici incepe ceva mult mai profund decat copiatul unei teme.

Profesorul isi poate pregati cursul cu AI. Elevul isi poate face tema cu
AI. Profesorul poate evalua cu AI. Elevul poate intreba AI-ul orice nu a
inteles si poate primi instantaneu un raspuns mai clar decat ar fi
reusit sa-si formuleze singur. Si, incet, fara sa existe o decizie clara
de a renunta la invatare, procesul prin care deveneai tu insuti mai
capabil poate fi inlocuit cu obtinerea rezultatului.

\textit{Nu mai inveti ca sa stii. Obtii raspunsul.}

\textit{Nu mai scrii ca sa formulezi o idee. Obtii textul.}

\textit{Nu mai stai sa intelegi de ce ai gresit. Intrebi ce trebuie
corectat.}

Si poate ca partea cea mai trista nu este nici macar asta.

Un prieten iti scrie ca ii este greu. Inainte, trebuia sa-l citesti, sa
te gandesti, sa incerci sa intelegi ce simte. Poate nu stiai ce sa-i
spui. Poate spuneai ceva nepotrivit. Poate il intrebai din nou. Poate
taceai cateva minute inainte sa raspunzi.

Acum poti intreba o masina ce sa-i raspunzi.

Si uneori ajungem la ceva si mai absurd: in loc sa-i raspundem noi, ii
trimitem direct conversatia cu AI-ul.

\textbf{Un om a venit la tine cu ceva al lui si tu ai transferat
intalnirea catre o masina.}

\textit{Cat de trist este sa ai in fata un om care incearca sa-ti spuna
ceva si sa-i raspunzi printr-o inteligenta care a inteles mesajul in
locul tau.}

Poate nici nu realizam toate aceste schimbari pentru ca AI-ul nu apare
intotdeauna sub forma unui lucru nou si spectaculos. Poate juca zeci de
roluri fara sa ne dam seama: profesor, coleg, secretar, terapeut,
editor, traducator, prieten, evaluator, adversar intr-un joc, partener
de conversatie. Fiecare rol, luat separat, pare util. Dar puse impreuna,
ele ating tot mai multe dintre locurile in care, pana acum, aveam nevoie
unii de altii.

Si atunci problema nu mai este doar daca AI-ul ne inlocuieste.

\textbf{Problema este daca incepe sa ne inlocuiasca relatiile cu noi
insine si cu ceilalti.}

Un student nu mai are nevoie sa-si construiasca rabdarea de a intelege.
Un profesor nu mai are nevoie sa-si pregateasca singur fiecare
explicatie. Un prieten nu mai trebuie ascultat pana la capat. Un om nu
mai trebuie sa gaseasca singur cuvintele pentru ceea ce simte.

Si, in paralel, apare comparatia.

Studentului i se spune ca trebuie sa fie mai productiv. Mai rapid. Mai
eficient. Apoi i se arata un sistem care citeste mai repede, scrie mai
repede, calculeaza mai repede, rezuma mai repede, raspunde mai repede.

\textit{``AI-ul face asta in zece secunde. Tie de ce ti-a luat doua
ore?''}

Pare o observatie despre productivitate.

\textbf{Dar repetata suficient de mult, devine o observatie despre
valoarea persoanei.}

\textbf{Nu mai compari doar rezultatul unei sarcini. Incepi sa compari
omul cu masina.}

Si comparatia este absurda din start. Nu compari doi oameni care au
primit aceleasi conditii. Compari un om - cu oboseala lui, cu
nesiguranta lui, cu timpul lui limitat, cu istoria lui - cu un sistem
construit tocmai pentru a executa anumite sarcini rapid si la scara.

Totusi, daca ii spui unui student suficient de des ca AI-ul face mai
bine ceea ce el incearca sa invete, apare intrebarea inevitabila:

\textit{``Atunci eu ce mai aduc?''}

Asta mi se pare una dintre cele mai periculoase parti ale schimbarii. Nu
faptul ca o masina poate face ceva mai bine decat mine, ci faptul ca as
putea incepe sa-mi masor propria valoare prin lucrurile pe care masina
le face mai bine.

\textbf{O sarcina poate fi inlocuita.}

\textbf{O persoana nu ar trebui redusa la suma sarcinilor pe care le
poate executa.}

Si poate ca aici apare, fara sa ne dam seama, o noua lovitura. Nu ni se
spune neaparat ca suntem mai putin importanti. Este suficient sa
incepem sa traim ca si cum am fi.

\textit{Sa nu mai invatam pentru ca AI-ul stie deja.}

\textit{Sa nu mai scriem pentru ca AI-ul o face mai bine.}

\textit{Sa nu mai ascultam pentru ca AI-ul stie ce sa raspunda.}

\textit{Sa nu mai incercam pentru ca undeva exista deja un sistem care
ar putea face lucrul acela mai repede.}

\textbf{Cea mai stranie forma de inlocuire nu este aceea in care AI-ul
ne ia locul. Este aceea in care noi ii facem loc renuntand singuri la al
nostru.} Inteligenta rezolva probleme. Constiinta inseamna sa existe ceva
ce se simte din interior in timp ce le rezolvi. Iar statutul moral
inseamna sa conteze pentru noi faptul ca acel ceva exista.

AI ne loveste deja foarte clar la prima intrebare. La celelalte doua nu
stim nici macar ce test am putea face.

Si poate asta ma tulbura mai mult decat posibilitatea ca AI sa fie mai
inteligent decat noi.

Pentru ca daca problema ar fi doar inteligenta, am putea accepta ca
exista ceva mai bun. Dar daca intr-o zi construim ceva care nu doar
vorbeste despre suferinta, ci chiar sufera, atunci nu mai avem o
problema de performanta.

\textbf{Avem o problema despre ce anume consideram ca are dreptul sa
conteze.}

Si aici se intoarce teatrul.

Ne place teatrul nu doar pentru ca vedem o poveste bine jucata, ci
pentru ca recunoastem pe scena ceva din noi. Iubirea, frica, rusinea,
empatia. Dar o parte din ceea ce ne misca vine si din faptul ca celalalt
putea sa rateze.

\textit{Actorul putea sa nu inteleaga.}

\textit{Putea sa uite.}

\textit{Putea sa ezite.}

\textit{Putea sa aleaga gresit.}

\textit{Si totusi sa gaseasca, uneori, cuvantul potrivit.}

O piesa care nu rateaza niciodata, care stie exact ce sa spuna fiecarui
spectator in fiecare secunda, ar putea fi mai eficienta decat teatrul.
Dar poate ca tocmai de aceea ar fi mai putin umana. \textbf{Ar semana
mai mult cu o forma perfecta de manipulare.}

La fel si intr-o scoala simulata. Un profesor artificial care stie
intotdeauna raspunsul, care intelege exact unde te-ai blocat si gaseste
de fiecare data formularea perfecta ar putea fi, in multe privinte, un
profesor mai bun decat mine.

Dar poate ca o parte din valoarea unui profesor nu sta in faptul ca este
de neinlocuit.

Poate sta si in faptul ca este un om care incearca, uneori fara sa stie
sigur ce face.

Poate ca sensul apare tocmai in locurile in care eficienta nu este
totul.

\textit{In esec.}

\textit{In ezitare.}

\textit{In faptul ca iti pasa de ceva care nu produce nimic.}

\textit{In faptul ca poti sa fii ranit de ceea ce face altcineva.}

\textit{In faptul ca poti sa nu stii ce sa spui si totusi sa ramai
acolo.}

\textbf{O viata eficientizata complet ar putea avea rezultate mai bune
si totusi sa nu mai stie de ce ar trebui sa le vrea.}

Si atunci revenim la lovitura de la inceput.

\textbf{Nu mai este doar cea pe care o citim in manuale. O traim deja,
zilnic.}

Deschidem telefonul si vedem miliarde de oameni care, la fel ca noi,
isi simt propria viata din centru. Fiecare are grijile lui, iubirile
lui, ambitiile lui, tragediile lui. Cu cat vezi mai multe centre
posibile, cu atat devine mai greu sa crezi ca sensul vine din simplul
fapt ca tu esti unul dintre ele.

Poate ca acesta este urmatorul pas.

Copernic ne-a spus ca nu suntem in centrul universului.

Darwin, ca nu suntem in centrul vietii.

Freud si neurostiinta, ca nu suntem nici macar in centrul propriei
minti.

Iar AI ne pune in fata unei intrebari si mai incomode:

\textit{Daca nici inteligenta nu este exclusiv a noastra, pe ce se mai
sprijina sentimentul ca noi suntem speciali?}

N-am de gand sa inchei cu un raspuns linistitor.

Poate ca maturitatea nu inseamna sa gasesti o cale de a te intoarce in
centru. Poate inseamna sa inveti sa construiesti sens fara sa ai nevoie
sa fii centrul.

Un copil are nevoie sa fie centrul universului cuiva.

Un adult poate sa predea o materie stiind ca, intr-o zi, un agent ar
putea sa o predea la fel de bine sau chiar mai bine.

Poate sa scrie stiind ca un model il poate imita.

Poate sa invete ceva ce o masina stie deja.

Poate sa raspunda unui prieten fara sa intrebe mai intai o inteligenta
artificiala ce ar trebui sa simta.

Nu pentru ca este de neinlocuit.

Ci pentru ca sensul actului nu a stat niciodata in faptul ca era
imposibil sa fie inlocuit.

\textbf{Predau cursul asta stiind exact lucrul asta.} Si totusi, poate
ca am pus prea mult accent pe momentul in care AI-ul va deveni suficient
de inteligent incat sa ne inlocuiasca. Pentru ca exista o forma mai
veche si mai discreta a aceleiasi probleme: algoritmii care nu incearca
sa faca lucrurile in locul nostru, ci sa decida ce vom vedea, ce ne va
atrage atentia si, incet, ce vom avea chef sa facem.

Deschizi TikTok-ul fara sa stii exact de ce. Un videoclip. Inca unul.
Inca unul. Nu ai intrat sa cauti ceva anume. Nu exista un punct in care
sa spui: ``gata, am terminat''. Feed-ul continua. Si, cu cat stai mai
mult, cu atat sistemul invata mai mult despre tine: ce ai urmarit pana
la capat, unde ai ezitat, ce ai revazut, ce ai ignorat, ce te-a facut sa
te opresti. Sistemele de recomandare sunt construite tocmai pentru a
personaliza cat mai bine ceea ce vezi si pentru a te mentine implicat;
in cazul TikTok, autoritatile europene au ridicat inclusiv problema
combinatiei dintre scroll infinit, autoplay, notificari si recomandari
foarte personalizate.

Problema nu este doar ca iti consuma timpul.

\textbf{Problema este ca iti consuma timpul fara sa-ti ceara sa decizi
ce vrei sa faci cu el.}

Poate aici apare o alta lovitura, mult mai subtila. Nu mai avem doar o
inteligenta care poate gandi in locul nostru. Avem sisteme care incearca
sa anticipeze ce vrem inainte sa stim noi ca vrem.

Nu mai alegi urmatorul videoclip. Sistemul il alege pentru tine. Nu mai
cauti urmatorul lucru care te intereseaza. Iti este oferit. Nu mai
trebuie nici macar sa formulezi o dorinta. \textbf{Trebuie doar sa
ramai.}

Iar asta este mult mai subtil decat sa-ti spuna cineva ce sa gandesti.
Pentru ca nu iti spune direct ce sa gandesti. Iti construieste mediul in
care vei gandi. Iti selecteaza urmatoarea imagine, urmatoarea poveste,
urmatoarea indignare, urmatoarea gluma, urmatorul om pe care il vei
compara cu tine.

\textbf{Si daca faci asta suficient de mult, incepe sa fie greu de spus
unde se termina preferinta ta si unde incepe ceea ce ai fost invatat sa
preferi.}

Poate ca lucrul cel mai tulburator este ca asta poate intra in viata
noastra mult mai adanc decat prin simplul consum de videoclipuri. Sa
spunem ca esti intr-o relatie. Ai un om langa tine si, ca in orice
relatie reala, apar momente in care te intrebi daca sunteti potriviti.
Aveti diferente. Te deranjeaza anumite lucruri. Poate esti ranit,
confuz sau pur si simplu treci printr-o perioada proasta si nu mai stii
exact ce simti.

Si atunci deschizi TikTok.

Algoritmul te ``citeste''. Vede la ce te opresti, ce urmaresti pana la
capat, ce revezi, ce cauti. Daca in perioada aceea incepi sa te uiti la
videoclipuri despre relatii, incompatibilitate, ``red flags'',
despartiri sau oameni care si-au dat seama prea tarziu ca erau cu
persoana nepotrivita, sistemul are toate motivele sa-ti ofere mai mult
din acelasi lucru. Nu pentru ca stie adevarul despre relatia ta. Nu
pentru ca a vorbit cu partenerul tau. Nu pentru ca a inteles ce se
intampla intre voi. Ci pentru ca a invatat ca \textbf{acel tip de
continut te face sa ramai.}

Si atunci se poate intampla ceva foarte ciudat. Tu ai pornit cu o
intrebare: ``Oare diferentele dintre noi sunt o problema?'' Iar
algoritmul incepe sa-ti ofere zeci, apoi sute de videoclipuri care par
sa raspunda: ``Da. Uite de ce.''

Nu neaparat pentru ca raspunsul este adevarat.

Ci pentru ca \textbf{indoiala ta produce retentie}.

Poate ca relatia chiar nu functioneaza. Poate ca videoclipurile acelea
pun degetul pe ceva real. Dar poate ca, la fel de bine, esti doar
obosit, nesigur, suparat sau intr-un moment dificil. Poate ca ceea ce ai
nevoie este sa vorbesti cu omul de langa tine. Poate ca ai nevoie de
timp. Poate ca ai nevoie sa descoperi singur daca problema este una
reala sau doar una trecatoare.

In schimb, primesti un flux continuu de oameni care iti explica propria
viata.

Si asta mi se pare trist.

Pentru ca algoritmul nu trebuie sa-ti spuna direct: ``Paraseste-ti
partenerul.'' Este suficient sa-ti ofere, timp de doua saptamani, sute
de oameni care iti explica de ce ar trebui.

Iar dupa suficiente videoclipuri, incepi sa ai impresia ca nu mai
privesti opiniile unor straini. Ai impresia ca ai descoperit un adevar
despre tine.

Poate ca algoritmul nu ti-a creat indoiala. Poate ca indoiala era acolo
de la inceput. Dar poate sa o hraneasca, sa o repete, sa o amplifice si
sa o transforme intr-o poveste coerenta inainte ca tu sa fi avut timp sa
stai singur cu ea.

Si aici devine important si ceea ce algoritmul nu iti arata.

Poate ca in aceeasi relatie exista lucruri frumoase. Poate ca omul de
langa tine te-a inteles intr-un moment in care nimeni altcineva nu a
facut-o. Poate ca radeti impreuna. Poate ca ati trecut prin lucruri
grele. Poate ca diferentele dintre voi nu sunt dovada ca nu sunteti
potriviti, ci pur si simplu diferente cu care inca invatati sa traiti.

Dar un videoclip despre o relatie sanatoasa, ambigua si complicata nu
produce neaparat aceeasi retentie ca unul care iti spune ca ``daca face
X, inseamna ca nu te iubeste''.

Algoritmul nu trebuie sa-ti controleze viata.

\textbf{Trebuie doar sa decida ce parte din ea vei vedea cel mai des.}

Si poate ca aici algoritmii ating aceeasi problema despre care vorbeam
cand discutam despre liberul arbitru. Libertatea nu inseamna doar sa
poti alege intre doua variante. Inseamna si sa ai sansa de a-ti forma
propria intrebare inainte ca altcineva sa inceapa sa-ti ofere
raspunsurile.

Sa poti avea o indoiala fara sa o transformi imediat intr-o concluzie.
Sa te poti plictisi. Sa stai intr-o camera fara sa ti se ofere imediat
ceva. Sa ai o conversatie dificila cu omul pe care il iubesti fara sa
cauti inainte zece straini care sa-ti spuna cine are dreptate. Sa te
razgandesti.

\textbf{Poate ca plictiseala este una dintre ultimele forme de spatiu
interior pe care le mai avem.}

Iar daca fiecare secunda goala este umpluta inainte sa apuci sa o simti,
pierdem ceva ce nu apare in nicio statistica de productivitate: timpul
in care nu faceam nimic si, tocmai de aceea, puteam deveni altcineva
decat ceea ce anticipase algoritmul.

De aceea, problema nu este doar ca AI-ul ar putea deveni mai inteligent
decat noi.

\textbf{Problema este ca, intre timp, construim sisteme din ce in ce mai
bune la a ne spune ce sa facem, ce sa privim, ce sa cumparam, cu cine sa
ne comparam si chiar cum sa ne interpretam propriile relatii.}

\textbf{Si toate acestea se intampla in timp ce avem senzatia ca suntem
perfect liberi.}

Poate ca urmatoarea mare lovitura nu va fi momentul in care o masina ne
spune ca este mai inteligenta decat noi.

Poate va fi momentul in care realizam ca, fara sa ne dam seama,
\textbf{am inceput sa vrem lucrurile pe care masinile au invatat sa ni
le ofere - si sa privim propria noastra viata prin povestile pe care ele
au invatat ca nu ne lasa sa plecam.}

\section*{Ce incerc sa insirui}

Un parcurs de la zero pana la nivel de olimpiada, care va evolua pe
parcurs. E facut din ce descriu ei pe ONIA, dar nu imi place ordinea in
care se face asta neaparat.

\begin{itemize}
  \item Python, NumPy, Pandas
  \item Prelucrarea si vizualizarea datelor
  \item Machine Learning fundamentals
  \item Regresie si clasificare
  \item Decision Trees, Random Forest, Gradient Boosting, SVM, KNN
  \item Clustering si PCA
  \item Feature engineering
  \item Train/validation/test si cross-validation
  \item Overfitting, underfitting si regularizare
  \item Metrici: Accuracy, Precision, Recall, F1, ROC-AUC, MSE, MAE, RMSE, R\textsuperscript{2}
  \item Hyperparameter tuning
  \item Neural Networks si PyTorch
  \item Computer Vision
  \item NLP
  \item Probleme si exercitii de tip olimpiada
\end{itemize}

Structura exacta nu este batuta in cuie. Cursul se va construi pe masura
ce vedem ce este util, ce este dificil si ce merita aprofundat.

\section*{Referinte}

\begin{itemize}
  \item \href{https://platform.olimpiada-ai.ro/ro/roadmap/programa}{Programa oficiala a Olimpiadei Nationale de Inteligenta Artificiala}
\end{itemize}
